% ---------------------------------------------------------------------------
% Appendix: Descriptions in a second domain — the menu-content split in
% deferred acceptance. Auctions are the paper's domain; this appendix replicates
% only the description finding (invariance content vs. mechanics restatement) in
% a non-auction strategy-proof mechanism, to show it is not auction-specific.
% The co-equal "second domain" material (full ranking, frontier grid, OSP
% protocol, Kendall-tau accounting) has been moved out of the paper.
% Sources: results/merged_ranking/da_cells.csv + da_cells_summary.md (menu-split
% cells; market-cluster bootstrap B=2000, numpy seed 1299; built by
% analysis/build_da_cells.py from the raw pipeline logs).
% ---------------------------------------------------------------------------
\section{Descriptions in a Second Domain: The Menu-Content Split in Deferred Acceptance}\label{app:da}

Auctions are this paper's domain. This appendix asks whether the description finding of Section~\ref{sec:ranking-descriptions}---that what moves behavior is not additional text about the mechanism but text that states the \emph{invariance} making truthful play safe---is specific to auctions, by replicating it in a second, non-auction strategy-proof environment: school choice under student-proposing deferred acceptance (DA). The setup is deliberately minimal. In each instance, four LLM students submit rank-order lists over four schools with unit capacity, under Ergin-acyclic priorities \citep{ergin2002efficient,ashlagi2018stable}; student-proposing DA is strategy-proof for the proposing side \citep{roth1982economics}, so a truthful list is dominant. We measure misreporting as the normalized Kendall-$\tau$ distance between the submitted and the true preference ranking (the fraction of the six school pairs ranked in opposite order; the DA analogue of the auction deviation metric of Section~\ref{sec:metrics}), and we run the same four model families as in the main text (Claude 3.5 Haiku, Gemini 2.0 Flash, GPT-4o, Gemma 27B). The pooled direct-DA baseline error is $4.2\%$. All numbers below are recomputed from the raw per-instance logs by \path{analysis/build_da_cells.py} (market-cluster percentile bootstrap, $B = 2000$, seed 1299); cross-model figures are unweighted means of per-model means.

\subsection{The Invariance Content Is the Active Ingredient}\label{app:da-descriptions}

The description cells (Family B; prompts in Appendix~\ref{app:prompts}) replicate, and sharpen, the auction-domain finding of Section~\ref{sec:ranking-descriptions}: what moves behavior is not additional text about the mechanism but text that states the \emph{invariance} making truthful play safe.

\paragraph{Safety-exposing description (B2).} The \emph{Rejection Safety} description states the property that distinguishes DA from the Boston mechanism \citep{abdulkadirouglu2003school}: rejections only redirect, never eliminate---a rejection from a school does not hurt the student's chances at her next-ranked school. It nearly eliminates misreporting, from $4.2\%$ to $0.2\%$ pooled (Claude $5.8\% \to 0.8\%$ $[0.3, 1.7]$; Gemini, GPT-4o, and Gemma all to exactly $0.0\%$).

\paragraph{Menu descriptions (B1) split by content.} \citet{gonczarowski2023strategyproofness} propose describing a strategy-proof mechanism through the \emph{menu} it induces: given others' reports, a student faces a fixed set of obtainable schools, and her report selects from that set. We implement two variants that bracket the informational content of this idea. The \emph{invariance-property} variant states the menu's incentive content explicitly (``your ranking \emph{cannot change} your obtainable set; it only determines which school you receive within it''). The \emph{mechanics} variant restates the same construction procedurally (``Step 1: the mechanism computes your obtainable schools. Step 2: you receive the school you ranked highest among them'') without asserting the invariance. The two variants move behavior in opposite directions:
\begin{itemize}
    \item \emph{Menu with the invariance property improves play}, behaving like a safety-exposing description: pooled error falls from $4.2\%$ to $1.8\%$ (Mann--Whitney $p < 10^{-4}$), with large significant gains for the two weak-baseline models (Claude $5.8\% \to 2.1\%$, MW $p = 10^{-4}$; Gemini $7.6\% \to 0.3\%$, $p < 10^{-4}$) and null effects for the models already near truthful (GPT-4o $p = 0.23$; Gemma $p = 0.43$).
    \item \emph{Menu mechanics alone worsens play}: pooled error rises from $4.2\%$ to $6.9\%$ ($p < 10^{-4}$), driven by GPT-4o ($1.0\% \to 4.2\%$) and especially Gemma ($2.6\% \to 13.7\%$). Restating the outcome computation gives the models material to strategize over without the property that would make strategizing pointless.
\end{itemize}
In auctions the menu restatement is null on average and sign-mixed across models (Section~\ref{sec:ranking-descriptions}); the DA data explain why: the standard menu text bundles a computation with an invariance, and only the invariance helps. Where the auction cell delivers the bundle and finds nothing, the DA cells unbundle it and find two opposing effects.

\begin{table}[htbp]
\centering
\footnotesize
\setlength{\tabcolsep}{4pt}
\begin{tabular}{llccccc}
\toprule
& & \multicolumn{4}{c}{Mean Kendall-$\tau$ error, \% [95\% CI]} & \\
\cmidrule(lr){3-6}
Description & Family & Claude & Gemini & GPT-4o & Gemma & Mean \\
\midrule
Direct DA (baseline) & --- & 5.8 [4.3, 7.5] & 7.6 [6.1, 9.2] & 1.0 [0.4, 1.7] & 2.6 [1.4, 4.0] & 4.2 \\
\addlinespace
Menu, mechanics only & B1 & 6.7 [5.3, 8.2] & 3.1 [2.0, 4.2] & 4.2 [2.9, 5.7] & 13.7 [11.0, 16.7] & 6.9 \\
Menu, invariance property & B1 & 2.1 [1.3, 3.0] & 0.3 [0.1, 0.7] & 0.3 [0.1, 0.7] & 4.2 [2.6, 6.1] & 1.8 \\
Rejection Safety & B2 & 0.8 [0.3, 1.7] & 0.0 [0.0, 0.0] & 0.0 [0.0, 0.0] & 0.0 [0.0, 0.0] & 0.2 \\
Textbook strategy-proofness & (advice) & 0.0 [0.0, 0.0] & 0.0 [0.0, 0.0] & 0.0 [0.0, 0.0] & 0.0 [0.0, 0.0] & 0.0 \\
\bottomrule
\end{tabular}
\caption{The menu-content split in DA: mean normalized Kendall-$\tau$ error in percent, with market-cluster bootstrap $95\%$ confidence intervals ($B = 2000$, seed 1299); ``Mean'' is the unweighted cross-model mean. Models: Claude 3.5 Haiku, Gemini 2.0 Flash, GPT-4o, Gemma 27B. Cells are $50$ markets $\times$ $4$ students (a few cells have $42$--$49$ markets after parse-failure drops), recorded in \texttt{results/merged\_ranking/da\_cells.csv}. A menu description that carries the invariance property improves play like a safety-exposing description; a menu description that restates only the mechanics worsens it. The textbook strategy-proofness cell states the dominant strategy outright and is included only as a diagnostic ceiling (see text).}
\label{tab:da-descriptions}
\end{table}

\paragraph{Diagnostic ceiling: the textbook statement.} A \emph{textbook strategy-proofness} description---a plain statement that submitting one's true ranking is optimal regardless of what others do---produces $0.0\%$ error for \emph{all four} models, marginally below even Rejection Safety ($0.2\%$). This cell reveals the dominant strategy outright and is not a simplicity-preserving lever, but it serves as a useful ceiling: when told the answer, every model complies perfectly, confirming that baseline misreporting reflects a failure to \emph{derive} strategy-proofness rather than an inability to execute a truthful report.

\paragraph{Removing the explanation.} Conversely, a \emph{null} description that states the rules with no explanation of the algorithm shifts play heterogeneously (pooled $4.2\% \to 5.2\%$, MW $p = 0.03$): it worsens Claude ($8.7\%$), GPT-4o ($3.3\%$), and Gemma ($7.0\%$) but \emph{improves} Gemini ($2.1\%$)---description content interacts with model-specific priors rather than acting uniformly.

\paragraph{Robustness: cardinal presentation.} An earlier variant of the grid presented raw dollar values (``$w = \$72$, $x = \$61$, \ldots'') instead of a sorted preference ordering, forcing the model to derive its own ranking. Baseline errors rise by an order of magnitude ($12$--$24\%$ across models), and the description split replicates: the invariance-property menu again produces the large improvement (Gemini $18.8\% \to 1.2\%$; GPT-4o $11.9\% \to 1.9\%$) while the mechanics variant barely moves the needle. The cardinal variant is incomplete (several cells were not run) and is documented, cell by cell, in \path{results/merged_ranking/da_cells.csv} (\texttt{variant=da\_cardinal}).
